% Configuration audit contribution to Part B. AI-produced by GPT-6 Astra.
% Provisional statement identifiers B.CF*; parent editor renumbers.
\subsection{Ordinary pro-reconstruction in finite transcendence degree}

We use an ordinary category here, not a category in which coherence is an
additional hypothesis. Let \(\mathsf T\) be the category whose objects are
Grothendieck topoi and whose morphisms are isomorphism classes of geometric
morphisms. Its ordinary pro-category has
\[
 \operatorname{Hom}_{\operatorname{Pro}(\mathsf T)}
       (\{\mathcal X_i\},\{\mathcal Y_j\})
   =\varprojlim_j\varinjlim_i
       \operatorname{Hom}_{\mathsf T}(\mathcal X_i,\mathcal Y_j).
                                                        \tag{B.CF.1}
\]
Fix a field \(k\) of characteristic zero and models of finite-type
\(k\)-schemes \(X,Y\) over a finitely generated subfield \(F_0\).
For finitely generated \(F\subset k\) containing \(F_0\), put
\[
 \mathcal X=\{X_{F,\mathrm{et}}\}_F,
 \qquad \mathcal Y=\{Y_{F,\mathrm{et}}\}_F,
 \qquad \mathcal B=\{\operatorname{Spec}(F)_{\mathrm{et}}\}_F.
\]
These are pro-objects of \(\mathsf T\), with structural arrows to
\(\mathcal B\). A morphism in the slice
\(\operatorname{Pro}(\mathsf T)_{/\mathcal B}\) induces a map
\(|X|\to |Y|\): apply the underlying-space functor and take inverse
limits. The identification with the scheme spaces follows from
\citep[Lemma~32.4.2, Tag~0CUF]{StacksProject}; its hypotheses are a directed
inverse system of schemes with affine transition maps, as here.
We call this pro-morphism \emph{pinned} when that induced map sends
closed points to closed points. This definition does not impose any
same-field representatives or coherence data.

\begin{claim}{B.CF1}{Proved}{Ordinary pro-reconstruction for every field of finite transcendence degree}
Suppose \(\operatorname{trdeg}_{\Q}k<\infty\). For finite-type
\(k\)-schemes, the functor to the preceding pro-systems is fully faithful
after universal homeomorphisms are inverted, with pinned morphisms in
\(\operatorname{Pro}(\mathsf T)_{/\mathcal B}\) as target morphisms.
Equivalently, for absolutely weakly normal finite-type \(X,Y\), every
such pinned ordinary pro-morphism is induced by a unique \(k\)-morphism
\(X\to Y\). This includes all finite-transcendence-degree fields of B8,
regardless of whether reconstruction from the single topos over \(k\)
holds.
\end{claim}
\begin{proof}
Enlarge \(F_0\) to contain a transcendence basis of \(k/\Q\).
Then \(k/F_0\) is algebraic. Every finitely generated subfield of \(k\)
is contained in one containing this \(F_0\), and every such field is
finite over \(F_0\). Thus we have made a cofinal change of the indexing
system, and all its field transitions are finite separable. In
particular \(X_E\) and \(Y_F\), for \(E\supset F\supset F_0\),
are finite-type schemes over the same finitely generated field \(F\).

We first verify that forgetting structural 2-cells causes no ambiguity
at these levels. For every finitely generated characteristic-zero
field \(F\), its absolute Galois group \(G_F\) is slim, meaning that
all its open subgroups have trivial centre. Indeed, if
\(\operatorname{trdeg}_{\Q}F>0\), apply
\citep[Corollary~4.5(1)]{HoshiTsujimura2023}; for a number field apply
Corollary~4.5(3) using any of its discrete mixed-characteristic valuation
rings. The same source, Definition~4.3, states the equivalent condition
that the centralizer in \(G_F\) of every normal open subgroup is trivial.
It follows also for every open subgroup, by taking its normal core.
For each connected component \(Z\) of a nonempty finite-type
\(F\)-scheme, the image of \(\pi_1(Z)\to G_F\) is open: it is the
stabilizer of a geometric connected component, of which there are
finitely many. An automorphism of the structural geometric morphism
\(Z_{\mathrm{et}}\to BG_F\) is an element of the centralizer of this
image. This follows by evaluating its action on every finite
\(G_F\)-set at a geometric base point; naturality and transport require
and suffice for commutation with the fundamental-group image.
Consequently that automorphism group is trivial. The same conclusion
holds component by component. Therefore, whenever the structural
triangle commutes as isomorphism classes, its commutation 2-cell is
unique. Moreover, any absolute isomorphism between two geometric
morphisms over \(BG_F\) automatically respects their structural 2-cells.

Let \(\alpha:\mathcal X\to\mathcal Y\) be a pinned pro-morphism over
\(\mathcal B\). At a target level \(F\), formula~(B.CF.1) represents
it by a geometric morphism
\[
             a:X_{E,\mathrm{et}}\longrightarrow Y_{F,\mathrm{et}}
                                                        \tag{B.CF.2}
\]
for some \(E\supset F\). The equality of structural pro-morphisms,
applied at the target level \(BG_F\), means that after enlarging \(E\)
once more, (B.CF.2) is a morphism over \(BG_F\). This is precisely the
equality criterion in the filtered colimit in (B.CF.1); it uses no
uniform choice over all target levels.

The representative (B.CF.2) is pinned. To prove this, let \(x\) be a
closed point of \(X_E\). The nonzero finite \(k\)-algebra
\(\kappa(x)\otimes_E k\) has a closed point, so \(x\) lifts to a
closed point \(\widetilde x\) of \(X\). The point
\(\widetilde y=|\alpha|(\widetilde x)\) is closed by assumption.
Its image in \(Y_F\) is closed: its residue field is generated by
finitely many coordinates over \(F\), is contained in a finite
extension of \(k\), and is algebraic over \(F\), since \(k/F\) is
algebraic. It is therefore finite over \(F\). Compatibility of the
underlying-space pro-morphism identifies that image with \(|a|(x)\).
This proves pinnedness at the actual finite stage.

We now apply the opened theorem
\citep[Theorem~0.1 and the equivalent formulation~(0.2)]{CarlsonHaineWolf2024}:
for a field finitely generated over its prime field, with the additional
hypothesis of positive transcendence degree only in positive
characteristic, the functor from finite-type schemes localized at
universal homeomorphisms to pinned relative geometric morphisms is
fully faithful. The field here is \(F\), of characteristic zero;
both schemes in (B.CF.2) are finite type over it. It follows that
(B.CF.2) comes from a unique localized \(F\)-scheme morphism
\(g_F:X_E\to Y_F\). Pairing this morphism with the structural map
to \(\operatorname{Spec}E\) gives a localized \(E\)-scheme morphism
\(X_E\to Y_E\), and then a localized \(k\)-scheme morphism
\(f_F:X\to Y\).

Choose this construction first at \(F_0\), and write \(f=f_{F_0}\).
For another target field \(F\), compatibility of the pro-morphism
means that \(g_F\) followed by \(Y_F\to Y_{F_0}\), and \(g_{F_0}\),
induce isomorphic geometric morphisms after a common finite source
extension \(H\subset k\). They are morphisms of finite-type
\(F_0\)-schemes. The uniqueness of structural 2-cells just proved
makes this a relative isomorphism; the faithfulness in Theorem~0.1
therefore makes the two localized scheme morphisms equal.
Their lifts to \(Y_F=Y_{F_0}\times_{F_0}F\) agree as well, since they
also have the same structural morphism to \(\operatorname{Spec}F\).
It follows that \(f_F=f\), and the pro-morphism induced by \(f\)
is exactly \(\alpha\), at every target level.

Conversely a localized \(k\)-scheme morphism descends to some finite
extension of \(F_0\): a representative finite diagram, its equations,
and its inverse universal homeomorphisms involve only finitely many
coefficients. Being a universal homeomorphism descends under the
faithfully flat field extension, and any necessary enlargement is
still finite over \(F_0\). Thus it defines the required ordinary
pro-morphism, and its map on limits is pinned. If two such morphisms
induce the same pro-morphism, compare them at target \(F_0\), refine
once, and apply the same finite-stage faithfulness and base change.
They are equal. This proves full faithfulness in the localized
category. The final formulation uses the equivalence of that
localization with absolutely weakly normal schemes, as in
\citep[Theorem~1.13, Lemma~1.32, Proposition~2.28]{CarlsonHaineWolf2024}.
\end{proof}

\begin{claim}{B.CF2}{Proved}{The remaining ordinary pro-category issue is an infinite-transcendence issue}
The mechanism in B.CF1 does not extend by the same argument to an
arbitrary characteristic-zero field: transitions corresponding to
transcendental field extensions need not preserve closed points.
Specifically, for \(F\) of characteristic zero and \(u\) transcendental,
the point \(T=u\) of \(\mathbb A^1_{F(u)}\) is closed, whereas its
image in \(\mathbb A^1_F\) is generic. Consequently the full model
system is not in general an object of the ordinary pro-category of
pinned geometric morphisms. This is an obstruction to that categorical
formulation, and is not a counterexample to reconstruction for pinned
morphisms in \(\operatorname{Pro}(\mathsf T)_{/\mathcal B}\).
\end{claim}
\begin{proof}
The contraction of the ideal \((T-u)\) from \(F(u)[T]\) to \(F[T]\)
is zero, because \(u\) is transcendental. This proves the assertion
on points. The scheme projection induces the transition geometric
morphism of the model system and has that same underlying-space map.
An inverse system in a category must have all its transition arrows
in that category. Finally, no assertion in this example constructs
a pinned ordinary pro-morphism failing to come from a scheme map;
thus it supplies no such counterexample.
\end{proof}

\begin{claim}{B.CF3}{Proved}{The infinite-transcendence obstacle and its resolution}
For \(\operatorname{trdeg}_{\Q}k=\infty\), ordinary relative
pro-reconstruction is proved by B.PRO1--5 below, culminating in
B.PRO5. The shorter argument of B.CF1 alone does not prove it:
a representative may have the form
\(X_{E,\mathrm{et}}\to Y_{F,\mathrm{et}}\) with \(E/F\)
transcendental. Pinnedness of the map on limits does not, by the
argument of B.CF1, give pinnedness of a finite-stage representative.
The coherent same-field result AS-FF.13 supplies neither that
finite-stage assertion nor an equivalence with this ordinary
pro-category. The new proof resolves the obstacle by deriving
coherence from compactness, comparing whole Kummer systems after
one finitely generated refinement, and reconstructing every étale
slice over a fixed source field. It does not assert that those
finite-stage transition morphisms preserve closed points.
\end{claim}

\subsection{A henselian field with no centralizing coefficient automorphisms}

\begin{claim}{B.CF4}{Proved}{A formal arc with no nontrivial algebraic self-correspondence}
There exists \(u\in t\Q[[t]]\), transcendental over \(\Q(t)\), with
the following property. Let \(\Omega\) be any complete rank-one valued extension
of \(\Q((t))\), with the valuation trivial on \(\Q\), and let
\(h\in\Omega\) satisfy \(v(h)=v(t)=1\). If both \(h\) and
\(u(h)\) are algebraic over \(\Q(t,u(t))\), and \(h,u(h)\) are
algebraically independent over \(\Q\), then
\[
                         h=t.                            \tag{B.CF.3}
\]
The series \(u\) can be specified by a recursive sequence of rational
finite jets; the proof gives a finite elimination certificate at every
stage. The property holds simultaneously in all finite ramified
extensions and their completions.
\end{claim}
\begin{proof}
Let \(Z\subset\mathbb A^2_{\Q}\times\mathbb A^2_{\Q}\) be an
irreducible closed subvariety of dimension two, dominant under both
projections, and different from the diagonal. We first prove that
in every cylinder of prescribed rational coefficients in \(t\Q[[t]]\)
there is a smaller cylinder none of whose series has a pair
\[
                 (t,u(t),h,u(h))\in Z(\Omega),\qquad v(h)=1.
                                                        \tag{B.CF.4}
\]
Write the prescribed jet as \(J(t)\) modulo \(t^N\), \(N\geq1\).
Consider the two-parameter family
\[
                       P(t)=J(t)+a t^N+b t^{N+1}.
\]
A generic member has finite intersection
\[
                         Z\cap(\Gamma_P\times\Gamma_P), 
                                                        \tag{B.CF.5}
\]
where \(\Gamma_P=\{(x,P(x))\}\). Here is a dimension proof, including
the loci where evaluation parameters coincide. On
\(xx'(x-x')\ne0\), evaluation at \(x,x'\) is a surjective linear
map from the two parameters \((a,b)\) to two values; the determinant
is \(x^N(x')^N(x'-x)\). The incidence variety of (B.CF.5) consequently
has dimension at most two over this locus. When \(x=x'\ne0\), the
point must lie in \(Z\cap\mathrm{Diag}\), of dimension at most one;
one independent evaluation condition again leaves total dimension at
most two. When \(x=0\), \(x'\ne0\), the first point is the fixed
point \((0,J(0))\). A fibre of the first dominant projection of the
irreducible surface \(Z\) has dimension at most one, and evaluation
at \(x'\) imposes one independent parameter condition. The same
argument applies with the factors exchanged. When both \(x=x'=0\),
only the single pair of fixed points is possible. Thus the full
incidence variety has dimension at most the two-dimensional parameter
space. Generic fibres are zero-dimensional or empty. We may choose
\(a,b\in\Q\), because rational points are Zariski dense in that
parameter space, including after passage to an algebraic closure.
This proves the assertion for the entire affine intersection; no
localization or denominator is needed.

Choose polynomial generators \(Q_j(x,y,x',y')\) of the ideal of
\(Z\). Since (B.CF.5) is zero-dimensional, its coordinate algebra is
finite-dimensional over \(\Q\). Elimination therefore supplies a
nonzero polynomial \(R(x)\in\Q[x]\) and polynomials
\(A_j(x,x')\in\Q[x,x']\) with
\[
 R(x)=\sum_j A_j(x,x')Q_j\bigl(x,P(x),x',P(x')\bigr).
                                                        \tag{B.CF.6}
\]
If \(u\equiv P\pmod{t^M}\), \(v(h)=1\), and (B.CF.4) holds,
then each summand on the right after \((x,x')=(t,h)\) has valuation
at least \(M\). Indeed \(u(t)-P(t)\) and \(u(h)-P(h)\) have
valuation at least \(M\), all rational coefficients have valuation
zero, and all remaining arguments have nonnegative valuation.
Taking \(M>\operatorname{ord}_{t=0}R(t)\), as well as
\(M>\deg P\), contradicts (B.CF.6). The same finite coefficient
cylinder is excluded for every valued extension \(\Omega\), not
just for one choice of Puiseux series field.

There are countably many such rational varieties \(Z\). Each
forbidden set is nowhere dense by the cylinder argument. The series
algebraic over \(\Q(t)\) also form a countable union of nowhere dense
sets: for every nonzero \(A(t,U)\in\Q[t,U]\), a polynomial \(P\)
with any prescribed jet can be chosen so that \(A(t,P(t))\ne0\),
and a sufficiently long further jet then preserves its nonzero
leading coefficient. The complete metric space \(t\Q[[t]]\), with
the coefficient topology, is a Baire space. Choose \(u\) outside
all these forbidden sets.

For the claimed pair \(h,u(h)\), take the rational Zariski closure of
\((t,u(t),h,u(h))\). Algebraicity and the two algebraic-independence
conditions make this an irreducible surface dominant on both factors.
If it were different from the diagonal it would be among the excluded
varieties. It is therefore the diagonal, which gives (B.CF.3).

For a recursive specification, enumerate rational polynomial ideals
which define the indicated irreducible dominant surfaces, and
interleave all nonzero polynomials \(A(t,U)\). At each step choose
the first rational pair \((a,b)\) in a fixed enumeration for which
(B.CF.5) is zero-dimensional, compute an elimination polynomial in
(B.CF.6), and extend the fixed jet beyond its order. Increase the
length strictly at every step. Polynomial ideal dimension,
irreducibility and elimination over \(\Q\) are finite algorithms;
the preceding proof guarantees that every search terminates. The
nested jets define a single series and prove its property. No finite
initial segment alone is asserted to establish all the infinitely
many exclusions.
\end{proof}

\begin{claim}{B.CF5}{Proved}{Trivial centralizers at every finite extension of a transcendence-degree-two henselian field}
Choose a series \(u\) as in B.CF4, and let
\[
             K=\Q(t,u(t))^h\subset\Q((t))               \tag{B.CF.7}
\]
be the henselization for the induced \(t\)-adic valuation. Then
\(K\) is countable, has transcendence degree two over \(\Q\), value
group \(\Z\), residue field \(\Q\), and completion \(\Q((t))\).
For every finite extension \(E/K\), in a fixed algebraic closure
\(\overline K\), one has
\[
              C_{\operatorname{Aut}(\overline K)}(G_E)=1.
                                                        \tag{B.CF.8}
\]
Thus the centralizer hypothesis of B7 is attained among the
residue-characteristic-zero fields with positive transcendence defect
specified by B8.
\end{claim}
\begin{proof}
The series is transcendental, so the initial field has transcendence
degree two. It lies in \(\Q((t))\), contains \(\Q(t)\), and has
residue field \(\Q\) and value group \(\Z\). Henselization is
algebraic and immediate. These properties, including its universal
property as a valued-field extension, are given in the opened
\citep[Proposition~3.3.25 and Corollary~3.3.26]{ADH2017}.
The polynomial truncations show that \(\Q(t)\) is dense in
\(\Q((t))\); hence the completion of \(K\) is that same field.
Countability follows from algebraicity over a countable field.

We record explicitly the group calculation used below. Inside a fixed
valued algebraic closure, put
\[
                  K^\dagger=\bigcup_{n\geq1}
                         K\overline{\Q}(t^{1/n}),
                                                        \tag{B.CF.9}
\]
with compatible roots. It is henselian, has algebraically closed
residue field \(\overline\Q\), and has divisible value group
\(\Q\). By the opened
\citep[Corollary~3.3.22]{ADH2017}, a valued field of equicharacteristic
zero is algebraically closed exactly when it is henselian with
algebraically closed residue field and divisible value group.
Thus (B.CF.9) is \(\overline K\). The constants \(\Q\) are
algebraically closed in \(K\), since \(K\subset\Q((t))\).
Adjoining a number field changes only the residue field, and adjoining
\(t^{1/n}\) then has degree \(n\), by its distinct valuation classes.
These facts show that restriction to constants and action on the
compatible roots give
\[
        G_K\simeq\widehat\Z(1)\rtimes G_\Q,             \tag{B.CF.10}
\]
where a splitting acts on the constants and fixes the chosen roots
of \(t\). The action on inertia is cyclotomic.

The centralizer in \(G_K\) of every open subgroup \(H\) is trivial.
An element of that centralizer projects to the centralizer of an
open subgroup of \(G_\Q\), which is trivial by
\citep[Corollary~4.5(3) and Definition~4.3]{HoshiTsujimura2023}.
It consequently lies in \(\widehat\Z(1)\). In each \(\ell\)-adic
coordinate, an element fixed by an open subgroup of \(G_\Q\) is
zero: the \(\ell\)-adic cyclotomic character is nontrivial on the
absolute Galois group of every number field, because a number field
contains only finitely many roots of unity, and multiplication by a
nonzero element of \(\Z_\ell\) is injective on \(\Z_\ell\).
This proves the asserted centralizer calculation without using the
conclusion (B.CF.8).

Now let \(E/K\) be finite and let
\(\sigma\in\operatorname{Aut}(\overline K)\) centralize \(G_E\).
Its fixed-field action preserves \(E\). The unique prolongation of
our valuation to \(E\) is henselian and discrete, with residue field
a number field, so \(E\) is not separably closed. The opened
\citep[Lemma~2.1]{KhassaKhanduja2010} states that every automorphism
of a non-separably-closed henselian valued field is continuous for
its valuation topology; its proof applies the theorem of F.\ K.
Schmidt to the original and transported valuations. For a discrete
rank-one valuation, equivalent valuation topologies give the same
valuation ring. An automorphism of its ordered value group is the
identity. Hence \(v(\sigma a)=v(a)\) on \(E\), and uniqueness of
prolongation over a henselian field gives the same equality on
\(\overline K\).

The field \(\overline\Q\subset\overline K\) is invariant under
\(\sigma\). The restriction of \(G_E\) to it is the open group
\(G_{E\cap\overline\Q}\); this follows from the Galois
correspondence, and the intersection is a number field since
\(E/K\) is finite and \(K\cap\overline\Q=\Q\). Therefore
\(\sigma|_{\overline\Q}\in G_\Q\) centralizes an open subgroup
of \(G_\Q\), and is the identity. In particular all coefficients
of \(u\) are fixed.

Embed \(\overline K\) as a valued field in an algebraic closure of
\(\Q((t))\), and work in its completion. Continuity and convergence
of the rational polynomial truncations give
\[
          \sigma u(t)=u(\sigma t),\qquad v(\sigma t)=1.
                                                        \tag{B.CF.11}
\]
Both \(\sigma t\) and \(\sigma u\) are algebraic over
\(\Q(t,u)\), because \(\overline K\) is algebraic over that
field; they are algebraically independent over \(\Q\), because
\(\sigma\) is a field automorphism. B.CF4 applied to
\(h=\sigma t\) yields \(\sigma t=t\), and (B.CF.11) gives
\(\sigma u=u\). The uniqueness of a valued embedding of a
henselization now makes \(\sigma\) the identity on \(K\).
It lies in \(G_K\) and centralizes its open subgroup \(G_E\).
The calculation following (B.CF.10) forces \(\sigma=1\), proving
(B.CF.8) for every finite \(E\).
\end{proof}

\subsection{A residue-characteristic-zero obstruction beyond coefficient automorphisms}

\begin{claim}{B.CF6}{Proved}{Compatible logarithmic reduction under a supported infinitesimal scaling}
Let \(R\) be a henselian discrete valuation ring of residue
characteristic zero, with fraction field \(K\), uniformizer \(t\),
and the unique extension \(v\) of its valuation to \(\overline K\).
Put \(\lambda=1+t\), and define
\[
 F(a)=
 \begin{cases}
   \lambda a,&v(a)>0,\\
   a,&v(a)\leq0,
 \end{cases}
 \qquad F(0)=0,
 \qquad F(\infty)=\infty.                                \tag{B.CF.12}
\]
For every finite \(G_K\)-stable set \(S\subset\mathbb P^1(\overline K)\)
containing \(0,1,\infty\), the stable marked genus-zero models for
\(S\) and \(F(S)\), over a finite Galois extension splitting their
marks, have isomorphic logarithmic special fibres over the same
standard log point. These isomorphisms preserve the labels and are
compatible with forgetting marks, finite extensions of the field of
definition, and the Galois descent data. At a mark \(a\), the induced
comparison of the arithmetic cusp group preserves its inertia and
induces the identity on its quotient \(G_{K(a)}\), with
\(K(F(a))=K(a)\) as actual subfields of \(\overline K\).
\end{claim}
\begin{proof}
The map (B.CF.12) is a \(G_K\)-equivariant bijection, whose inverse
scales the same open valuation ball by \(\lambda^{-1}\).
For two distinct finite points,
\[
          \frac{F(a)-F(b)}{a-b}\in1+\mathfrak m_{\overline K}.
                                                        \tag{B.CF.13}
\]
When both points belong to the ball this ratio is \(\lambda\), and
when neither does it is one. In the remaining case, say
\(v(a)>0\geq v(b)\), the perturbation is \(ta\), which has
valuation strictly larger than \(v(a-b)=v(b)\). This proves
(B.CF.13). In particular every difference has the same valuation and
the same initial form in the associated graded valued field before
and after applying \(F\). The assertion also holds for differences
in the reciprocal coordinate near infinity: \(F\) fixes all points
in that branch.

Choose a finite Galois extension \(L/K\) containing all the marks,
and write \(\pi\) for a uniformizer of its valuation ring. Distinct
rational sections of \(\mathbb P^1_{R_L}\) can be separated by
successive blowups at their common reductions. Each blowup is at a
smooth point of the special fibre. Its exceptional component has an
affine coordinate
\[
                          z=(x-a)/\pi^m,                 \tag{B.CF.14}
\]
for a centre \(a\) in the corresponding cluster of marks, with
\(m\) the depth of that cluster. Equivalently one may use the
reciprocal coordinate on the infinity branch. The finite tree of
components, their depths, and the locations of the marked and nodal
points on each component are obtained from the valuations and
normalized initial forms of the differences between the marks.
Contracting components with fewer than three special points gives
the stable marked model. This description is also a direct
construction of it over \(R_L\).

Equation (B.CF.13) identifies the two trees, their depths, and all
component configurations. In the coordinates (B.CF.14), the
component isomorphisms have the same reduced coordinates. They are
unique on stable components because three labelled special points
fix an automorphism of \(\mathbb P^1\). Equip the models with the
fs divisorial log structures defined by the special fibre and the
marked sections. A node on the stable model has a toric chart
\(uv=\epsilon\pi^d\), with \(d\geq1\) and \(\epsilon\) a
unit. The logarithmic special fibre retains \(d\) and the residue
of the smoothing unit, with its two tangent directions. In cluster
coordinates before contraction one may write
\(uv=\pi^{n-m}\); after changing normalized coordinates the
smoothing unit is a product and quotient of initial forms of the
same differences. Thus (B.CF.13) preserves the smoothing data and
the component maps lift to log isomorphisms over the fixed log point.
These toric charts show at the same time that the total marked model
is log regular; contracting a chain merely replaces thickness one
by a positive integral thickness.

We give the compatibility calculation, since uniqueness of the
isomorphisms by itself would not suffice for contractions. If
\(z=(x-a)/\pi^m\) is a coordinate for a target component and
\(w=(x-b)/\pi^n\) is one for a source component, the transition is
\[
       z=\frac{b-a}{\pi^m}+\pi^{n-m}w.                   \tag{B.CF.15}
\]
Extract the smaller valuation of the two coefficients. Its first
nonzero normalized reduction is a nonzero constant, \(w\), or
\(c+w\); it depends only on the initial form of \(b-a\) and the
two depths. The transformed expression has exactly the same
exponent and the same reduced factor, by (B.CF.13). This checks
both the characteristic monoid and the unit part of the pullback
of each logarithmic generator, including a generator which
vanishes identically on a contracted tail. All node, mark, and
infinity charts are obtained by products, quotients, and
reciprocals of these linear expressions. Hence the logarithmic
isomorphisms commute with every mark-forgetting contraction.
Equality of the unit factors may be checked on the irreducible
components, because the special fibre is reduced.

For completeness, the lift of a fixed component isomorphism is
unique once the characteristic monoids and log base are fixed.
At a node its possible change multiplies the two generators by
units \(u_1,u_2\) with \(u_1u_2=1\), \(u_1x=x\), and
\(u_2y=y\) in the reduced nodal ring. Thus
\(u_1-1\in(x)\cap(y)=0\). On the smooth locus the vertical
generator comes from the fixed log base, and a horizontal generator
maps to a non-zero-divisor. There is no additional freedom there.
The same coordinate check as in (B.CF.15) works under finite base
extension: if \(\pi=\epsilon'(\pi')^e\), the same unit
\(\epsilon'\) occurs on both sides, and initial forms commute
with the map between the associated graded valued fields.
This proves base-extension compatibility. Since (B.CF.12) is
\(G_K\)-equivariant, the uniquely labelled component maps and their
log lifts commute with \(\operatorname{Gal}(L/K)\).

At a marked section (a), the local horizontal parameter is
\(x-a\). Its image is \(x-F(a)\); inside the ball their
comparison is the scaling \(\lambda\), and outside it the
identity. The vertical log structure and the field \(K(a)\)
are unchanged, since multiplication by the nonzero element
\(\lambda\in K\) preserves \(K(a)\). The local Kummer chart
therefore identifies the horizontal inertia \(\widehat\Z(1)\)
and its arithmetic quotient \(G_{K(a)}\), with the identity on
the latter. This is the usual cusp description obtained from the
punctured strict henselian normal trait, or directly from the
Kummer chart in SGA~1, XIII, Proposition~5.2. Choices of geometric
fibre functors change the resulting group maps by the corresponding
cuspidal conjugacy, not by an unspecified automorphism of the
residue field.
\end{proof}

\begin{claim}{B.CF7}{Refuted}{Reconstruction over an equicharacteristic-zero henselian discretely valued field}
For every field \(K\) admitting a nontrivial henselian discrete
valuation with residue characteristic zero, closed-point-preserving
reconstruction of finite-type \(K\)-schemes from their relative
small étale topoi is false, even with universal homeomorphisms
inverted. The supported scaling (B.CF.12) is the geometric-point
function of a pinned autoequivalence
\[
          (\mathbb A^1_K)_{\mathrm{et}}
             \longrightarrow(\mathbb A^1_K)_{\mathrm{et}}
                  \quad\hbox{over }BG_K                 \tag{B.CF.16}
\]
which is not induced by a morphism of \(K\)-schemes.
\end{claim}
\begin{proof}
We first state precisely the opened comparison theorem.
\citep[Theorem~1.6]{Lepage2013} assumes that \(\mathscr X\) is a
log-regular fs log scheme and every prime in \(L\) is invertible
on it. Restriction identifies its finite Kummer-étale covers of
pro-\(L\) type with the corresponding finite étale covers of its
log-trivial locus. Theorem~1.8 in the same paper assumes a
noetherian henselian base \(B\) with closed point \(s\), a
connected fs log scheme \(\mathscr X\), and properness of its
underlying scheme over \(B\); restriction identifies finite
Kummer-étale covers of \(\mathscr X\) with those of its log
special fibre. These are equivalences of covering categories,
functorial in the log schemes. For our models all primes are
invertible: their valuation rings contain \(\Q\). Taking
\(L\) to be the set of all primes gives, for each marked model,
\[
 \operatorname{FEt}(\mathbb P^1_L\setminus S)
       \simeq\operatorname{KCov}(\mathscr X_S)
       \simeq\operatorname{KCov}((\mathscr X_S)_s).
                                                        \tag{B.CF.17}
\]
Here \(L/K\) now denotes a finite splitting field for the marks;
the prime-set letter in the quoted theorem has served its purpose.
The model is proper over a henselian DVR and log regular by the
charts in B.CF6. Thus every hypothesis has been checked. This
application retains all finite covers, with no prime omitted.

Apply (B.CF.17) to the logarithmic special-fibre isomorphisms of
B.CF6. We obtain compatible inverse-image equivalences
\[
 \operatorname{FEt}(\mathbb P^1_L\setminus F(S))
           \longrightarrow
 \operatorname{FEt}(\mathbb P^1_L\setminus S).            \tag{B.CF.18}
\]
Their compatibility with mark-forgetting, finite constant
extensions, and Galois descent was checked in (B.CF.15), so
finite étale descent gives the same system over \(K\).
These equivalences are over \(\operatorname{FEt}(K)\), rather
than merely geometric: apply the two quoted theorems also to
\(\operatorname{Spec}R\) with its vertical log structure.
They identify covers of the fixed log point with finite étale
\(K\)-algebras, and every special-fibre isomorphism used above
lies over that same log point. At each puncture (B.CF.18)
identifies the full decomposition and inertia pair, and acts as
the identity on the actual quotient \(G_{K(a)}\), as proved
at the end of B.CF6.

We now pass from this compatible system of finite covering
categories to the entire small étale site. This step is necessary;
an isomorphism at one punctured fundamental group alone would
not suffice. A quasi-compact separated étale object
\(V\to\mathbb P^1_K\) is a disjoint union of normal curves.
Zariski's Main Theorem identifies it with an open subscheme of
the finite normal compactification \(\overline V\to\mathbb P^1_K\)
in its generic finite separable algebras. Its complement is a
finite set of points. Choose a finite set \(S\) containing the
branch locus, the images of every point of
\(\overline V\setminus V\), and \(0,1,\infty\).
Over its complement the compactification is finite étale.
Use (B.CF.18) to transport this
finite étale cover, and normalize \(\mathbb P^1_K\) in the
resulting generic algebra. The full cusp pairs give the
bijections between the points over the omitted base points:
ramification is measured by inertia, and the unramified residue
extension by the quotient decomposition group over the fixed
\(G_{K(a)}\). Therefore they transport the removed points of
\(\overline V\setminus V\), preserve the unramified locus, and
produce a quasi-compact separated étale object again.
The construction is independent, up to its compatible natural
isomorphisms, of enlarging \(S\), by the restriction
compatibility in (B.CF.18).

A morphism between such étale objects is determined by its map
of generic finite algebras, extended uniquely to their normal
compactifications, together with the condition that its image
lies in the chosen open subsets. The covering-category
comparisons preserve these generic maps and the conditions on
points. Thus they give an equivalence of these categories of
étale objects. They preserve covering families: joint
surjectivity at a geometric generic point is detected by the
finite covering category; at a geometric closed point it is
detected by the same local decomposition/inertia comparison.
These objects form a generating small étale site. We have
therefore constructed an autoequivalence of
\((\mathbb P^1_K)_{\mathrm{et}}\), over \(BG_K\), whose
function on \(\overline K\)-valued points is exactly (B.CF.12). It is pinned,
since that function is a \(G_K\)-equivariant bijection and
preserves the finite field \(K(a)\) at every algebraic point.
It fixes infinity and hence restricts to (B.CF.16).

Every nonzero rational number is a valuation unit and is fixed
by (B.CF.12), whereas \(t\) is sent to \((1+t)t\ne t\).
If a \(K\)-scheme morphism \(\mathbb A^1_K\to\mathbb A^1_K\)
induced (B.CF.16), its polynomial would agree with \(T\) on the
infinite subset \(\Q^\times\), and would consequently be \(T\)
identically. Its value at \(t\) gives the contradiction.
The scheme \(\mathbb A^1_K\) is already absolutely weakly
normal, so inverting universal homeomorphisms introduces no
additional morphism which could realize this point function.
\end{proof}

\begin{claim}{B.CF8}{Refuted}{The centralizer condition in B7 is insufficient}
There exists a characteristic-zero field \(K\), countable and of
transcendence degree two over \(\Q\), such that
\[
 C_{\operatorname{Aut}(\overline K)}(G_E)=1
           \quad\hbox{for every finite }E/K,
\]
but closed-point-preserving reconstruction of finite-type
\(K\)-schemes from their relative étale topoi fails. One may take
exactly the recursively specified field (B.CF.7). The obstruction
is a compatible nonalgebraic permutation of points and their
whole covering categories, not a centralizing field automorphism.
\end{claim}
\begin{proof}
B.CF5 proves the displayed condition for this field, with all
finite extensions included. Its valuation is henselian and
discrete, with residue field \(\Q\), so B.CF7 supplies the
pinned nonalgebraic autoequivalence of its affine-line topos.
These are the same field and the same entire arithmetic base
\(BG_K\). This is a counterexample to the asserted implication.
We also verify the surviving necessary direction, rather than
accepting AS-FF.6 on authority. If a field satisfies reconstruction,
then every finite extension does: finite étale slicing identifies
objects and relative morphisms over that extension, and faithfulness
forces the reconstructed map to preserve the extension's structural
map. Let \(E/K\) be finite and let \(\sigma\) centralize
\(G_E\). It preserves \(E\). The coefficient automorphism on
\(\mathbb G_{m,E}\), together with the identification of its
base Galois action supplied by \(\sigma G_E\sigma^{-1}=G_E\)
pointwise, gives a pinned relative topos automorphism whose geometric
point function is \(a\mapsto\sigma(a)\). This includes the case
where \(\sigma\) fixes \(E\): then it is the identity absolute
morphism with its relative structural 2-cell twisted by the central
element. The point-function description follows directly by pulling
back to finite separable extension points, as in the opened
\citep[Proposition~3.6]{CarlsonHaineWolf2024}.
Reconstruction over \(E\) would make this an \(E\)-automorphism
of \(\mathbb G_m\), hence \(T\mapsto cT^{\pm1}\).
Evaluation at one gives \(c=1\); evaluation at two excludes
inversion. Consequently \(\sigma(a)=a\) for every
\(a\in\overline K^\times\), so \(\sigma=1\).
This proves necessity for all finite extensions. B.CF8 shows that
excluding this particular obstruction is not sufficient.
A precise valid sufficient-and-necessary replacement is the
point-function condition of the opened
\citep[Theorem~4.18(5)]{CarlsonHaineWolf2024}: for every affine
regular connected finite-type source with a rational point,
every pinned relative map to \(\mathbb G_m\) must have the
geometric-point function of a regular unit. The supported
scaling demonstrates why a condition solely on coefficient
field automorphisms cannot replace that requirement.
\end{proof}

\begin{claim}{B.CF9}{Refuted}{The specified transcendence-defect field in AS-FF.18 does not reconstruct}
Let
\[
     u_*(t)=\sum_{n\geq1}t^{n!},\qquad
     K_* =\Q(t,u_*(t))^h\subset\Q((t)).                  \tag{B.CF.19}
\]
This field has transcendence degree two, residue field \(\Q\),
value group \(\Z\), and positive transcendence defect one.
It does not satisfy closed-point-preserving étale reconstruction.
This settles the particular residue-characteristic-zero case
left undecided in AS-FF.18, without claiming that its centralizers
have the property proved for the different series in B.CF4.
\end{claim}
\begin{proof}
For a nonzero \(P(t,U)\in\Q[t,U]\), write \(e=\deg_U P\)
and \(m=\deg_t P\), and put
\(u_N=\sum_{n=1}^N t^{n!}\). If \(P(t,u_*)=0\), then a
nonzero \(P(t,u_N)\) has order at least \((N+1)!\) but degree
at most \(m+eN!\), impossible for sufficiently large \(N\).
Thus infinitely many distinct \(u_N\) would be roots of the
nonzero polynomial \(P(t,U)\) over \(\Q(t)\), again impossible.
The series is transcendental. The field and its henselization
therefore have the asserted transcendence degree, and
henselization preserves residue field and value group.
The defect is \(2-0-1=1\). B.CF7 applies directly to this
henselian discrete valuation and refutes reconstruction.
\end{proof}

\begin{claim}{B.CF10}{Open}{Exact remaining scope of the henselian-field question B8}
The universal quantifier in B8 is not fully settled. B.CF7--9,
together with B.LIT5--6, refute reconstruction for every field of
finite transcendence degree admitting a nontrivial henselian
valuation of residue characteristic zero. This includes all the
residue-characteristic-zero transcendence-defect cases in AS-FF.18,
the specified field \(K_*\), and fields with a mixed-characteristic
valuation admitting such a nontrivial residue-characteristic-zero
coarsening. The following mixed-characteristic cases remain from
the original \texttt{agent\_F/F\_section.tex}.
\begin{enumerate}
\item General finite-transcendence-degree henselian fields of mixed
residue characteristic outside the positive and negative classes
already separated in AS-FF.16, AS-FF.22, and AS-FF.26--39.
In particular, an unbounded wild ramification index is not by itself
a negative criterion: the cyclotomic, maximal abelian, and each fixed
bounded-nilpotency class in AS-FF.35, 37, 39 were already positive
claims in the return.
\item For every prime \(p\), the specified maximal metabelian field
\[
 K_p=\operatorname{Frac}(\Z_{(p)}^{sh}),\qquad
 I=G_{K_p},\qquad
 B_p=\overline\Q^{\,\overline{[I',I']}},\quad
 I'=\overline{[I,I]},
\]
and general finite extensions of it. AS-FF.40 singles out this
field; AS-FF.18 names \(B_2\) as the remaining degree-zero test.
Neither its asserted centralizer condition nor the affine and
selected-cover rigidity assertions AS-FF.44--50 prove reconstruction
for arbitrary pinned morphisms or arbitrary covers. The negative
residue-characteristic-zero argument here has no application to its
mixed-characteristic valuation.
\item AS-FF.39 explicitly makes no assertion about the union of
the fixed nilpotency-class fields
\(\bigcup_c\overline\Q^{\,\gamma_{c+1}(I)}\), or about a general
solvable closure. The former union remains undecided here. If the
latter phrase means the maximal solvable extension of this precise
strict-local field \(K_p\), B.CF11 below identifies it with
\(\overline\Q\) and refutes reconstruction. It does not identify
the lower-central union with the derived-series union, and makes no
assertion about solvable closures of arbitrary other base fields.
B.CF13 determines the stable rational Kummer kernel at the
lower-central union for \(p\geq5\), with explicit exceptions
at \(p=2,3\), but its full reconstruction problem remains open.
\end{enumerate}
\end{claim}
\begin{proof}[Scope verification]
The original pending list is stated at AS-FF.18 and its subsequent
updates: AS-FF.27 handles algebraic fields with a discrete
mixed-characteristic valuation; AS-FF.29--30 handle the specified
bounded-wild fields; AS-FF.35, 37, 39 handle the cyclotomic, maximal
abelian, and fixed bounded-nilpotency extensions. AS-FF.40--50 retain
the full reconstruction problem over \(B_p\), while proving only the
more restricted assertions indicated above. None of these scope
statements converts a failed proof mechanism into a counterexample.
Theorem B.CF7 uses a discrete residue-characteristic-zero valuation
for its noetherian stable logarithmic models. B.LIT5 removes
discreteness by the verified valuative-logarithmic approximation
and covering-category continuity; B.LIT6 removes the rank-one
restriction in finite transcendence degree by henselian coarsening.
These additions settle the entire residue-characteristic-zero
portion. They do not treat positive residue characteristic and
its wild covers. The full classification demanded by B8 therefore
remains open in the mixed-characteristic respects listed above.
\end{proof}

\begin{claim}{B.CF11}{Refuted}{The unbounded solvable closure of the strict-local field}
Fix an embedding of the strict henselization
\(A=\Z_{(p)}^{sh}\) in \(\overline\Q\), put
\(K_p=\operatorname{Frac}(A)\), and write \(I=G_{K_p}\).
Let \(I^{(0)}=I\) and
\(I^{(d+1)}=\overline{[I^{(d)},I^{(d)}]}\) be the closed
derived series. Then
\[
 K_p^{\mathrm{sol}}
   :=\bigcup_{d\geq 0}\overline\Q^{\,I^{(d)}}
   =\overline\Q.                                      \tag{B.CF.20}
\]
Consequently this maximal solvable closure does not satisfy
closed-point-preserving \(\acute etale\) reconstruction. The statement
concerns the actual field in AS-FF.39, not a substitution of a
completed local field. It leaves open the bounded derived-length
field \(B_p=\overline\Q^{I^{(2)}}\), and the lower-central union
\(\bigcup_c\overline\Q^{\gamma_{c+1}(I)}\).
\end{claim}
\begin{proof}
The ring \(A\) is a strictly henselian discrete valuation ring
with residue field \(\overline{\F}_p\). For a finite Galois
extension \(M/K_p\), henselianity gives a unique extension of
the valuation. Thus its decomposition group is the whole
\(\operatorname{Gal}(M/K_p)\). Its residue extension is trivial,
since it is finite over \(\overline{\F}_p\), so the whole
Galois group is also its inertia group. The opened
\citep[Lemma~15.114.5, Tag~09EE]{StacksProject} gives an exact
sequence with a finite \(p\)-group as kernel and a cyclic
prime-to-\(p\) group as quotient. Both are solvable, hence every
finite quotient of \(I\) is solvable. Equivalently one may apply
the profinite version, Lemma~58.13.5, Tag~0BUA, to this same DVR.

Every nonidentity element of a profinite group survives in some
finite quotient. If that quotient is solvable of derived length
\(d\), the image of \(I^{(d)}\) in it is trivial. It follows
that \(\bigcap_d I^{(d)}=1\). Alternatively, every element of
\(\overline\Q\) lies in a finite Galois extension of \(K_p\),
whose finite solvable Galois group has some finite derived length;
that extension lies in \(\overline\Q^{I^{(d)}}\). This proves
(B.CF.20), without asserting any uniform bound on derived length.

Choose a nonidentity automorphism \(\sigma\) of
\(\overline\Q/\Q\), for example complex conjugation after
choosing a complex embedding. Acting on coefficients induces an
autoequivalence of \(\mathbb A^1_{\overline\Q,\mathrm{et}}\)
over \(B G_{\overline\Q}=\mathsf{Sets}\); it preserves closed
points and acts on their values by \(a\mapsto\sigma(a)\).
If induced by a morphism of \(\overline\Q\)-schemes, its
polynomial would fix every rational number, and hence would be
\(T\). It moves an algebraic number, a contradiction. Inverting
universal homeomorphisms does not change this conclusion for the
normal characteristic-zero affine line. This is a genuine
negative reconstruction result for the exact unbounded derived
tower, rather than a failure of a proposed proof of reconstruction.
\end{proof}
